\documentclass[11pt]{article}
\usepackage{ap-physc}
\course{AP Physics C: Mechanics}
\headcourse{AP Physics C: Mechanics}
\headpart{Unit 2 \textperiodcentered\ Question Set}
\title{Unit 2: Force and Translational Dynamics: Question Set}
\author{Paul Wang}
\date{November 2026}
\docpart{Unit 2 \textperiodcentered\ Question Set}

\begin{document}
\maketitle

\begin{abstract}
Exam-style practice on Unit 2 (topics 2.1--2.10), which is 20--25\% of the multiple-choice section. Section~1--3 hold 13 four-choice multiple-choice questions grouped by topic; most need a short calculation or a line of reasoning, as on the exam. The last page before the key is a full Translation Between Representations free-response question (12 points). All items are original and written for practice. Use the AP equation sheet and a calculator. Allow about 26 minutes for the multiple choice (the exam's pace is about 2 minutes per question) and 24 minutes for the free response. The answer key explains every item and names its topic; a scoring guide for the free response follows it.
\end{abstract}

\scoringline{MCQ: \blank\ / \qtotal \qquad FRQ: \blank\ / 12 \qquad Target: 9+ MCQ \qquad Pace: 2 min per MCQ}

\section{Systems and Newton's laws (2.1--2.5)}

\begin{q}{3}{2.1 Center of mass}{$x_{cm} = (2.0 \cdot 0 + 6.0 \cdot 4.0)/8.0 = 3.0$ m, closer to the heavier mass.}
A 2.0 kg mass sits at $x = 0$ and a 6.0 kg mass at $x = 4.0$ m. Where is the center of mass of the pair?
\chh{$1.0$ m}{$2.0$ m}{$3.0$ m}{$4.0$ m}
\end{q}
\begin{q}{2}{2.1 Continuous center of mass}{$M = \int_0^2 3x^2\,dx = 8$ kg and $\int_0^2 x\cdot 3x^2\,dx = 12$ kg$\cdot$m, so $x_{cm} = 1.5$ m.}
A rod 2.0 m long lies from $x = 0$ to $x = 2.0$ m. Its linear density is $\lambda = 3x^2$ (kg/m). Where is its center of mass?
\chh{$2.0$ m}{$1.5$ m}{$1.3$ m}{$1.0$ m}
\end{q}
\begin{q}{1}{2.3 Third law}{Newton's third law: the two forces are a pair and are always equal and opposite. The car's larger acceleration comes from its smaller mass, $a = F/m$.}
A loaded truck collides head-on with a small car. During the collision, how does the force the truck exerts on the car compare with the force the car exerts on the truck?
\ch{They are equal in magnitude and opposite in direction.}{The truck exerts the larger force because it has more mass.}{The car exerts the larger force because it has the larger acceleration.}{The truck exerts the larger force because it is moving faster.}
\end{q}
\begin{q}{2}{2.5 Atwood machine}{First $a = (5.0 - 3.0)(9.8)/8.0 = 2.45\ \mathrm{m/s^2}$. Then for the 3.0 kg mass, which accelerates upward, $T = m_1(g + a) = 3.0(12.25) = 37$ N.}
Masses of 3.0 kg and 5.0 kg hang from the two ends of a light string over an ideal pulley and are released. What is the tension in the string? ($g = 9.8\ \mathrm{m/s^2}$)
\chh{$29$ N}{$37$ N}{$39$ N}{$49$ N}
\end{q}
\begin{q}{3}{2.5 Apparent weight}{Slowing while moving up means the acceleration is downward, $a = -2.0\ \mathrm{m/s^2}$. Then $N = m(g + a) = 60(7.8) = 468$ N.}
A 60 kg student stands on a scale in an elevator. The elevator is moving upward but slowing down at $2.0\ \mathrm{m/s^2}$. What does the scale read? ($g = 9.8\ \mathrm{m/s^2}$)
\chh{$710$ N}{$590$ N}{$470$ N}{$120$ N}
\end{q}
\begin{q}{2}{2.5 Time-varying force}{$a = F/m = 3t$, so $v = \int_0^2 3t\,dt = \tfrac32 t^2\big|_0^2 = 6.0$ m/s.}
A net force $F = 6t$ (newtons, $t$ in seconds) acts on a 2.0 kg block that starts from rest. What is its speed at $t = 2.0$ s?
\chh{$3.0$ m/s}{$6.0$ m/s}{$12$ m/s}{$24$ m/s}
\end{q}

\section{Gravity, friction, and springs (2.6--2.8)}

\begin{q}{4}{2.6 Gravitational field}{The distance from Earth's center doubles ($r = 2R_E$), and $g \propto 1/r^2$, so $g = 9.8/4 = 2.45\ \mathrm{m/s^2}$.}
A satellite orbits at an altitude equal to Earth's radius. What is the gravitational field strength at its location? ($g = 9.8\ \mathrm{m/s^2}$ at the surface)
\chh{$4.9\ \mathrm{m/s^2}$}{$9.8\ \mathrm{m/s^2}$}{$0\ \mathrm{m/s^2}$}{$2.5\ \mathrm{m/s^2}$}
\end{q}
\begin{q}{3}{2.7 Static friction}{The maximum static friction is $\mu_s mg = 19.6$ N, more than 15 N, so the crate stays put and static friction just matches the push.}
A 4.0 kg crate rests on a level floor ($\mu_s = 0.50$, $\mu_k = 0.30$). A worker pushes it horizontally with 15 N. What is the friction force on the crate? ($g = 9.8\ \mathrm{m/s^2}$)
\chh{$0$ N}{$12$ N}{$15$ N}{$20$ N}
\end{q}
\begin{q}{2}{2.7 Constant-velocity incline}{Constant velocity: $mg\sin\theta = \mu_k mg\cos\theta$, so $\mu_k = \tan 30^\circ = 0.58$.}
A block slides down a ramp inclined at $30^\circ$ at constant velocity. What is the coefficient of kinetic friction?
\chh{$0.50$}{$0.58$}{$0.87$}{$1.7$}
\end{q}
\begin{q}{4}{2.8 Springs in series}{Each spring carries the full weight $mg$ and stretches $mg/k$; two in series stretch $2mg/k$. Equivalently $k_{eq} = k/2$.}
Two identical springs, each of constant $k$, are hooked end to end and hung from the ceiling. A block of mass $m$ hangs from the bottom and is at rest. By how much does the pair stretch in total?
\chh{$mg/k$}{$mg/(2k)$}{$4mg/k$}{$2mg/k$}
\end{q}

\section{Drag and circular motion (2.9--2.10)}

\begin{q}{1}{2.9 Linear drag}{$v(t) = v_T(1 - e^{-t/\tau})$ with $\tau = m/b$. At $t = \tau$, $v = (1 - e^{-1})v_T = 0.63\,v_T$.}
A ball is dropped from rest in a fluid that exerts a drag force $-bv$. Its terminal speed is $v_T$. What is its speed at time $t = m/b$ after release?
\chh{$0.63\,v_T$}{$0.37\,v_T$}{$0.50\,v_T$}{$0.86\,v_T$}
\end{q}
\begin{q}{3}{2.10 Flat curve}{Static friction supplies the centripetal force: $\mu_s mg = mv^2/r$, so $v = \sqrt{\mu_s g r} = \sqrt{392} = 19.8$ m/s.}
A car rounds a flat curve of radius 50 m. The coefficient of static friction between tires and road is 0.80. What is the greatest speed at which it can take the curve without skidding? ($g = 9.8\ \mathrm{m/s^2}$)
\chh{$28$ m/s}{$22$ m/s}{$20$ m/s}{$14$ m/s}
\end{q}
\begin{q}{4}{2.10 Vertical circle}{At the bottom $T - mg = mv^2/r$, so $T = mg + mv^2/r$; at the top $T = mv^2/r - mg$.}
A ball on a string is whirled in a vertical circle at nearly constant speed. Where is the string's tension greatest?
\ch{At the top, where gravity and tension both point toward the center}{At the sides, where gravity is perpendicular to the string}{It is the same everywhere because the speed is constant}{At the bottom, where tension must both balance gravity and supply the centripetal force}
\end{q}

\section{Free response}

\begin{frq}{Question 1 \textperiodcentered\ Translation Between Representations \textperiodcentered\ 12 points \textperiodcentered\ 24 minutes}
A small sphere of mass $m$ is released from rest at $t = 0$ in a tall column of oil. The oil exerts a drag force of magnitude $bv$, where $b$ is a positive constant and $v$ is the sphere's speed. Ignore buoyancy. Take downward as positive.

\begin{enumerate}[label=(\alph*),leftmargin=2em,itemsep=4pt]
\item Draw and label a free-body diagram of the sphere at an instant after release when it is moving downward. Represent each force with an arrow starting on the dot. \pts{2}
\item Write, but do not solve, a differential equation for $v(t)$. Then solve it, using the initial condition, to derive $v(t)$ in terms of $m$, $b$, $g$, and $t$. \pts{4}
\item Sketch a graph of $v$ versus $t$. On the graph, label the terminal speed as an expression and indicate the slope of the graph at $t = 0$. \pts{3}
\item A second sphere of the same size, so the same $b$, but mass $2m$ is released in the same oil. On the same axes, how does its graph compare with the first? Justify your answer using your expressions from part (b). \pts{3}
\end{enumerate}
\end{frq}

\makekey

\subsection*{Free-response scoring guide (12 points)}
\small
Award each point only when the answer does what the model does. Points are independent unless a part says otherwise; a correct method with an answer carried forward from an earlier error earns the later points.
\begin{scoring}
(a) (2) & 1 pt: a downward arrow labeled weight or $mg$. 1 pt: an upward arrow labeled drag, $bv$, with no extra forces (no ``net force'' or velocity arrows). & Adding a separate ``terminal force'' arrow, or drawing drag downward. \\
(b) (4) & 1 pt: $m\,\dfrac{dv}{dt} = mg - bv$. 1 pt: separates variables, $\displaystyle\int_0^v \frac{dv'}{g - (b/m)v'} = \int_0^t dt'$. 1 pt: integrates correctly, $-\frac{m}{b}\ln\!\left(1 - \frac{bv}{mg}\right) = t$. 1 pt: $v(t) = \dfrac{mg}{b}\left(1 - e^{-bt/m}\right)$. & Writing $v = gt - (b/m)vt$, which treats $v$ as constant inside the drag term. \\
(c) (3) & 1 pt: curve starts at the origin and rises, concave down. 1 pt: approaches a horizontal asymptote labeled $v_T = mg/b$ without crossing it. 1 pt: initial slope labeled $g$ (drag is zero at $v = 0$). & A straight line up to $v_T$ that then turns flat, or a curve that overshoots $v_T$. \\
(d) (3) & 1 pt: same initial slope $g$ (both start with only gravity acting). 1 pt: terminal speed doubles, $v_T = 2mg/b$. 1 pt: it takes longer to approach terminal speed because $\tau = 2m/b$ doubles; its curve lies above the first at every $t > 0$. & Claiming the heavier sphere has a larger initial acceleration. \\
\end{scoring}

\end{document}
